% !TeX root = RJwrapper.tex
\title{Stabilitest: Robustness and Fragility Analysis of Statistical Test Conclusions in R}


\author{by Marius Ardelean}

\maketitle

\abstract{%
A p-value quantifies how surprising the observed data would be under the null hypothesis, but not how easily the resulting decision could be overturned. The stabilitest package combines three sensitivity views of one pre-specified analysis: jackknife leave-one-out influence, greedy worst-case observation removal, and the bootstrap same-decision rate. Component metrics and a composite 0-100 score are returned for two-sample location and proportion tests, linear model and ANCOVA terms, GLM and Cox terms, and TOST equivalence and non-inferiority endpoints. Categorical interpretation is fail-closed: a label is emitted only where an independent pre-registered study validates thresholds for that exact configuration. A Fisher exact-test study validated a two-band mapping. A prospective Welch study completed 4,500 training analyses but found no feasible mapping, so held-out data remained unopened and Welch labels are suppressed. The article also reports negative ANCOVA calibration studies. These outcomes illustrate why numeric sensitivity diagnostics and categorical assurance claims require different evidentiary standards.
}

\section{Introduction}\label{introduction}

Statistical reporting in clinical trials typically compresses the evidence into
a binary decision at \(p < \alpha\). That compression discards something a
reader needs. Consider two trials that both report \(p = 0.048\). In the first,
the conclusion survives essentially any plausible perturbation of the data. In
the second, it rests entirely on one extreme responder: delete that single
patient and significance evaporates. Standard reporting treats these as
identical findings.

This is not a hypothetical concern: sensitivity of conclusions to individual
patients is a recurring theme in the trial-reporting literature
\citep{senn2007, pocock2015}. The ICH E9 guideline and its E9(R1) addendum
\citep{ich1998e9, ich2019e9r1} make it explicit, asking that inferences be robust
to data limitations and to deviations from the assumptions of the main
estimator, assessed through pre-specified sensitivity analyses. What the
guidance does not supply is quantitative machinery. Three distinct questions jointly characterise the
stability of a test conclusion, and each needs a different tool:

\begin{enumerate}
\def\labelenumi{\arabic{enumi}.}
\tightlist
\item
  \emph{Which individual observations drive the result?} (influence)
\item
  \emph{What is the smallest set of observations whose removal would overturn the
  conclusion?} (fragility)
\item
  \emph{How often do bootstrap resamples preserve the observed decision?}
  (same-decision rate)
\end{enumerate}

\CRANpkg{stabilitest} \citep{ally2026stabilitest} answers all three for one
pre-specified primary analysis, returning component metrics plus a composite
score. Its contributions are:

\begin{itemize}
\tightlist
\item
  \textbf{An integrated framework.} Jackknife influence, adversarial worst-case
  removal, and the bootstrap same-decision rate are computed together for a common
  catalogue of trial-relevant analyses --- two-sample location and proportion
  tests, linear model and ANCOVA terms, GLM and Cox terms, and TOST
  equivalence and non-inferiority --- through a small, uniform API.
\item
  \textbf{A fail-closed calibration registry.} Categorical labels
  (\emph{Robust} / \emph{Fragile}) are treated as claims that require evidence. The
  package emits a label only where an independent, pre-registered calibration
  study has validated thresholds for that exact method, endpoint, conclusion
  type, weight vector, and design profile. Everywhere else the label is \texttt{NA}
  while numeric output is retained. Thresholds are never transferred between
  method families by analogy.
\item
  \textbf{A calibrated unit, earned.} An independent study for Fisher's exact test
  yielded a validated two-band mapping, reported here with its held-out
  confirmation bounds.
\item
  \textbf{Named negative results.} A prospective Welch study and three ANCOVA
  studies closed without usable thresholds under their frozen designs. These
  empirical failures are reported rather than converted into weaker claims.
\end{itemize}

Deletion sensitivity is complementary to the estimand and sensitivity-analysis
strategy required by ICH E9(R1). It identifies dependence on observed rows; it
does not by itself address missing-data assumptions, intercurrent events, or
alternative estimands.

\section{Related R packages}\label{related-r-packages}

Every package in Table \ref{tab:related} answers a version of ``how much would
have to change for this conclusion to flip?'', but they fall into two families.

\CRANpkg{sensemakr} \citep{cinelli2020, cinelli2024sensemakr},
\CRANpkg{konfound} \citep{rosenberg2025konfound} and \CRANpkg{EValue}
\citep{vanderweele2017evalue}
ask the question about \emph{unmeasured} threats --- a confounder, a selection
mechanism, or unobserved cases the analyst never had. Their answers are
analytic bounds computed from the reported estimate and its uncertainty; none
of them touches the rows of the analysis dataset. \CRANpkg{boot}
\citep{canty2025boot} is a general resampling engine rather than a conclusion-level
robustness framework. \CRANpkg{fragility}
\citep{lin2025fragility}, \CRANpkg{influence.ME} \citep{nieuwenhuis2012influenceme},
\CRANpkg{car} \citep{fox2019car}, \texttt{zaminfluence} (GitHub only; \citet{giordano2026zaminfluence},
with the associated paper \citet{broderick2023}) and \CRANpkg{stabilitest} instead ask about the
\emph{observed} sample: which specific patients, if perturbed or removed, would
change the conclusion?

Within that second family the differences are of scope and rigour rather than
kind. \CRANpkg{fragility} is restricted to flipping binary outcome labels, implementing
the clinical fragility index of \citet{walsh2014}. \CRANpkg{car} and \CRANpkg{influence.ME}
offer classical single-deletion diagnostics in the tradition of \citet{cook1977},
which answer a different, single-deletion question from the package's
multi-observation adversarial search. \texttt{zaminfluence}
approximates the adversarial leave-\(k\)-out answer efficiently via influence
functions, but for Z-estimators in general rather than the two-sample and
model-term designs of trial reporting.

\begin{table}

\caption{\label{tab:related}R packages addressing related sensitivity questions. Detailed endpoint and method comparisons are supplied in related-packages.md. Entries were checked against CRAN or repository documentation on 2026-08-10; fragilityindex is excluded because it has been archived on CRAN since 2020-07-08.}
\centering
\begin{tabular}[t]{lll}
\toprule
Package & Threat & Scope\\
\midrule
stabilitest & Observed-row deletion & Tests and model terms\\
fragility & Outcome-status flips & Binary outcomes\\
sensemakr & Omitted confounding & Linear regression\\
konfound & Omitted bias / cases & Regression inference\\
EValue & Unmeasured bias & Observational effects\\
\addlinespace
boot & Sampling variation & Any statistic\\
influence.ME & Influential groups & Mixed models\\
car & Single-case influence & lm / glm\\
zaminfluence & Approximate deletion & Z-estimators\\
\bottomrule
\end{tabular}
\end{table}

\CRANpkg{stabilitest}'s distinguishing choices relative to all of these are: it
evaluates deletion, jackknife and bootstrap sensitivity together as three
views of one conclusion rather than offering a single diagnostic; its
worst-case removal is an actual adversarial refit search rather than a
residual/leverage heuristic or a linear approximation; and its composite score
carries a categorical verdict only where an independent pre-registered
calibration has validated it for that exact configuration. No other package in
Table \ref{tab:related} ships a calibration registry or a fail-closed
labelling policy.

\section{Methods}\label{methods}

Given two samples --- or a model term, see the API section --- and a
significance level \(\alpha\), four analyses are performed. For two independent
groups, \(n_1\) and \(n_2\) denote per-group sizes and \(N=n_1+n_2\) is the total
number of observations. Three analyses enter the composite score; the fourth
is retained as a deliberate contrast.

\textbf{Jackknife (leave-one-out).} Each observation is deleted in turn and the
test repeated, following the classical framework of \citet{shaotu1995}. Outputs are
the conclusion stability \(S_\mathrm{jack}\) (percentage of deletions preserving
the significance decision), the set of influential observations (deletion
flips that decision, or shifts \(p\) by more than \(\delta = 0.05\)), and the leave-one-out
\(p\)-value range. Because single-deletion influence shrinks as \(\sim 1/n\), high
jackknife stability in a large sample is \emph{weak} evidence of robustness. The
component's chief value is identifying specific observations for data-quality
review; the package's printed output says so explicitly. Conclusion stability
does not guarantee preservation of effect direction: direction reversals must
be inspected separately from the significance decision.

\textbf{Worst-case removal (primary fragility analysis).} A greedy algorithm
deletes, at each step, the observation whose removal moves the \(p\)-value
furthest toward overturning the conclusion, up to
\(k_{\max} = \lfloor 0.30\,n \rfloor\) deletions. The \emph{removal fragility index}
\(k_\mathrm{frag}\) is the step at which the conclusion first flips. If no flip
occurs by \(k_{\max}\), the stored value \(k_{\max}+1\) is a right-censoring code
and is reported as \(>k_{\max}\), not as an observed flip. This
follows the maximum influence perturbation of \citet{broderick2023} (see also
\citet{giordano2026a} and \citet{giordano2026b}), who showed that in prominent published
analyses removing under 1\% of observations can reverse the sign of an
estimated effect. When a flip is found, the greedy set is \emph{sufficient} to
overturn the conclusion, so \(k_\mathrm{frag}\) is an upper bound on the size of
the minimal overturning subset; no universal optimality claim is made. Because
a flip from \(p = 0.047\) to \(p = 0.053\) is a knife-edge
event \citep{potter2020}, the package returns and plots the full \(p\)-value
trajectory rather than the index alone. Note the terminology: the \emph{removal}
fragility index counts deleted observations and so reduces \(n\), which
distinguishes it from the Walsh fragility index \citep{walsh2014}, which flips
event status while holding \(n\) fixed.

\textbf{Extreme-value removal (descriptive).} Cumulative deletion of observations
most distant from the recomputed grand mean --- the familiar ``we removed the
outliers and the result held'' procedure. It does \emph{not} enter the composite
score. It is retained because the gap between this index and the worst-case
index measures how much protection the outlier narrative actually provides.

\textbf{Bootstrap same-decision rate.} \(B\) resamples are drawn with replacement
within groups \citep{efron1993}; \(S_\mathrm{boot}\) is the empirical proportion that
preserves the observed significance decision. It is a plug-in quantity under
the empirical distribution, not an unconditional probability that a future
study will replicate. It mostly reflects strength of evidence at the observed
effect size, not resistance to contamination: a marginal \(p\)-value can yield
\(S_\mathrm{boot}\) near 50--60\% in clean data. The API retains the historical
field name \texttt{bootstrap\_reproducibility} for compatibility.

\textbf{Composite score.}
\(R = w_1 S_\mathrm{jack} + w_2 F + w_3 S_\mathrm{boot}\), where
\(F = 100 \cdot \min\!\left(k_\mathrm{frag} / (k_{\max}+1),\, 1\right)\)
rescales the fragility index to the full 0--100 range. Default weights are
\(w = (0.4, 0.4, 0.2)\): influence and fragility carry equal weight, the
same-decision term is down-weighted because it largely restates the
\(p\)-value. The weights are a heuristic communication device, not a derived
quantity. They are an explicit argument and should be fixed in the statistical
analysis plan. Two results with the same score can be stable in different
ways, so the components --- especially the fragility index and the identity of
the observations behind it --- are the substance of reporting.

\section{Package design and API}\label{package-design-and-api}

\subsection{Dispatcher and engines}\label{dispatcher-and-engines}

Two-sample analyses go through a single dispatcher; model-based and
equivalence analyses have their own entry points sharing a common
case-deletion engine.

\begin{verbatim}
# Two-sample: t.test, paired.t.test, wilcoxon, wilcoxon.paired,
#             brunner_munzel, fisher, chisq, prop
robustness_analysis(group1, group2, test_type = "t.test")

# Model terms: single coefficient or a multi-df factor tested jointly
robustness_lm(change ~ arm + baseline, data = d, term = "arm")
robustness_glm(y ~ arm + age, data = d, term = "arm", family = binomial())
robustness_surv(Surv(time, event) ~ arm, data = d, term = "arm")

# Equivalence / non-inferiority: mean, risk difference, odds ratio
robustness_tost(g1, g2, type = "equivalence", endpoint = "mean",
                delta_L = -20, delta_U = 20)
\end{verbatim}

For model engines all components operate on \emph{rows} (subjects) of the analysis
dataset, with the conclusion defined by the \(p\)-value of a pre-specified term
--- a single coefficient, or a multi-df factor term via a joint test
(\texttt{drop1} \(F\) for \texttt{lm}, likelihood-ratio for \texttt{glm} and \texttt{coxph}). Deleting whole
subjects, with their covariates and follow-up, preserves the adjustment
structure and handles censoring naturally. Fits that fail after deletion
(a factor level vanishing, say) are skipped and reported rather than silently
dropped.

Arguments are deliberately uniform across engines: \texttt{alpha}, \texttt{n\_boot},
\texttt{max\_removal\_pct}, \texttt{weights}, and \texttt{seed} mean the same thing everywhere, and
\texttt{seed} defaults to \texttt{123} so that a bare call is reproducible.

\subsection{Return objects and methods}\label{return-objects-and-methods}

Analyses return a classed list (\texttt{robustness\_analysis}, \texttt{robustness\_model}, or
\texttt{robustness\_tost}) containing the original test result, a
\texttt{robustness\_metrics} tibble of 17 component metrics, the per-step trajectories
of each removal analysis, the bootstrap draws, the resolved calibration
decision, and the analysis profile used to make that decision. \texttt{print()}
methods render a structured summary; \texttt{plot()} returns the removal trajectories
and the bootstrap \(p\)-value distribution.

\subsection{The calibration registry}\label{the-calibration-registry}

The mechanism that makes label suppression systematic rather than ad hoc is a
registry shipped as \texttt{inst/extdata/calibration-registry.csv}, keyed by the
\emph{resolved} method rather than by the user-facing argument. A call with
\texttt{test\_type\ =\ "t.test"} resolves to \texttt{welch\_unpaired} or \texttt{paired\_t}; one with
\texttt{test\_type\ =\ "fisher"} resolves to \texttt{fisher\_exact}. Each row carries a status,
cutoffs where applicable, a version string, and a provenance pointer to the
study that produced it.

Label resolution is fail-closed. A categorical label is emitted only when the
resolved method has a \texttt{validated\_method\_specific} row \textbf{and} the observed
conclusion type matches \textbf{and} the weight vector matches the calibrated one
\textbf{and} the design profile falls inside the study's eligibility bounds. Any
mismatch yields \texttt{NA} --- with a printed reason --- while the numeric score and
all components are retained. There is no fallback to a ``nearest'' calibration.
The historical Welch 55/70 mapping is now an inactive comparator and is never
applied to Welch, ANCOVA, or another runtime result.

\section{Calibration policy and evidence}\label{calibration-policy-and-evidence}

\subsection{The Welch unit}\label{the-welch-unit}

The original 55/70 mapping was based on differences between scenario-level
average scores. A prospective replacement study tested whether those averages
supported case-level categorical interpretation. The calibration unit was a
significant independent-groups Welch test at \(\alpha=0.05\), with default
weights \((0.4,0.4,0.2)\), \texttt{n\_boot\ =\ 1000}, and a 30\% removal cap. The six core
design archetypes crossed per-group sizes from 25 to 200, balanced and 40/60
allocations, and variance ratios 1, 2, and 1/2. Within each archetype the null
used zero mean difference, the borderline stratum targeted 60\% power, and the
clear stratum targeted 95\% power.

All 18 training cells reached their quota of 250 completed significant
analyses: 1,500 null, 1,500 borderline, and 1,500 clear analyses (4,500 total),
with zero analysis failures. The frozen search considered integer three-band
and two-band mappings. No candidate passed all safety, identification,
ordinal-accuracy, occupancy, and ordering requirements, so the final result is
\texttt{uncalibrated} with no cutoffs. Under the pre-specified protocol this training
outcome was final: held-out validation remained unopened, no validation
estimates exist, and no positive runtime eligibility profile was published.

The historical 55/70 mapping was retained as a named training comparator. Its
false-reassurance estimate was 0.458
(Wilson upper 0.479,
scenario-cluster upper 0.513,
conservative upper 0.513),
against gates of 0.05 and 0.10. Clear-effect identification was
0.499
(Wilson lower 0.478,
scenario-cluster lower 0.348,
conservative lower 0.348),
against gates of 0.70 and 0.60. Balanced ordinal accuracy was
0.558 (gate 0.70),
and minimum band occupancy was
0.227. Pooled ordering
was preserved and no material reversal was detected, but those secondary
conditions could not rescue the failed primary gates.

The frozen no-candidate SHA-256 is represented in eight consecutive blocks.
The first four are \texttt{9c45481b} \texttt{952cab7c} \texttt{b9b9086e} \texttt{37924a39}; the final four
are \texttt{d83fe048} \texttt{4628745d} \texttt{be2e79eb} \texttt{33e8797d}.
the compact supplementary artifact records this hash, source hashes, runtime,
occupancy, and the unopened-validation flag. The active registry version is
\texttt{welch-2026-2}: numeric scores and components remain available, while the
categorical label is \texttt{NA}.

Table \ref{tab:simtable} retains the earlier 12-scenario experiment as a
descriptive stress example, regenerated with the released package (12
scenarios \(\times\) 500 replications, \(B=200\)). It was not used to fit or
confirm the prospective decision.

\begin{table}

\caption{\label{tab:simtable}Descriptive historical Welch stress simulation, conditional on statistical significance; not training or confirmation evidence for the prospective calibration decision. There are 500 replications per scenario and B = 200. Columns report effect size, per-group n, injected contaminating observations, rejection rate, mean score among significant results, median worst-case and extreme-value fragility indices, and mean bootstrap same-decision rate.}
\centering
\begin{tabular}[t]{rrrrrrrr}
\toprule
$d$ & $n$ & Cont. & Rej. rate & Score & $k_{\mathrm{wc}}$ & $k_{\mathrm{ex}}$ & Boot. \%\\
\midrule
0.0 & 25 & 0 & 0.040 & 54.7 & 2.0 & 2 & 67.1\\
0.0 & 25 & 2 & 0.064 & 52.6 & 1.0 & 1 & 66.6\\
0.0 & 50 & 0 & 0.076 & 50.0 & 1.5 & 2 & 63.1\\
0.0 & 50 & 2 & 0.084 & 51.3 & 2.0 & 2 & 65.5\\
0.5 & 25 & 0 & 0.392 & 61.1 & 3.0 & 4 & 75.3\\
\addlinespace
0.5 & 25 & 2 & 0.670 & 64.5 & 4.0 & 4 & 80.2\\
0.5 & 50 & 0 & 0.698 & 63.2 & 5.0 & 11 & 80.3\\
0.5 & 50 & 2 & 0.864 & 65.7 & 6.0 & 10 & 84.8\\
0.8 & 25 & 0 & 0.808 & 69.7 & 5.0 & 8 & 83.0\\
0.8 & 25 & 2 & 0.922 & 73.3 & 6.0 & 9 & 88.6\\
\addlinespace
0.8 & 50 & 0 & 0.976 & 75.5 & 12.0 & 26 & 94.1\\
0.8 & 50 & 2 & 1.000 & 77.6 & 14.0 & 26 & 96.3\\
\bottomrule
\end{tabular}
\end{table}

Two features of Table \ref{tab:simtable} deserve emphasis. First, under the
null the rejection rate is at or near nominal in clean data
(0.040 and 0.076 across the two
sample sizes) and is higher under contamination
(0.064 and 0.084). The Monte Carlo
standard error at 500 replications is about 0.010, so individual
cells should not be over-read. Among findings that reached significance,
median worst-case fragility is one or two observations throughout the null
rows. This is a descriptive pattern, not evidence that a particular score
threshold safely classifies an individual result.

Second, and more consequentially, the median extreme-value index is
systematically larger than the worst-case index in the \(d=0.8\) scenarios ---
26 versus 12 at \(d = 0.8\), \(n = 50\).
Judging stability by removing apparent outliers makes a result look roughly
twice as sturdy as an adversarial reading shows it to be. The most damaging
observations are usually not the obvious ones.

Figure \ref{fig:simfig} shows the underlying score distributions.

\begin{figure}
\includegraphics[width=1\linewidth,alt={Two faceted histograms of Welch robustness scores. Null and large-effect score distributions overlap, and dashed vertical lines mark the inactive historical cutoffs at 55 and 70.}]{stabilitest_files/figure-latex/simfig-1} \caption{Descriptive score distributions for significant results under the null (d = 0, left) and the study's d = 0.8 scenario (right), n = 50 per group, uncontaminated. Dashed lines mark the inactive historical 55/70 mapping. The distributions overlap; this figure is not validating evidence for categorical bands.}\label{fig:simfig}
\end{figure}

\subsection{The Fisher unit: a calibration earned}\label{the-fisher-unit-a-calibration-earned}

An independent study for Fisher's exact test on two-arm binary outcomes
delivered the package's sole active categorical calibration. Its design was frozen
before execution: truth strata defined by \emph{enumerated exact power} rather than
by observed effect size (null exact; clear at exact power 0.95; a borderline
stratum at 0.60 retained as diagnostic only and excluded from threshold
fitting), and a sealed score-only pilot gating production.

The study also settled the weight question empirically. On binary data the
jackknife stability component is constant at 100 in every cell --- a single
deletion cannot move a Fisher \(p\)-value across \(\alpha\) at these sample sizes
--- so it carries zero information. The frozen design therefore uses
\emph{jackknife-light} weights \((0, 0.5, 0.5)\), which maximise score granularity
without changing discrimination.

Training selected the cutoff \(L = 58\) with false reassurance 0.0493 (Wilson
upper 0.0607) and identification of clear effects 0.7608 (Wilson lower 0.740).
Held-out confirmation, opened only once, gave 0.0367 and 0.7333; under the
conservative-of-Wilson-and-cluster-bootstrap rule the reported bounds are a
false-reassurance upper of 0.0617 and an identification lower of 0.6767, both
inside the frozen gates. The registry row is \texttt{fisher-2026-1}, provenance
\texttt{study:binary\_proportion@cc3344931614}.

The resulting mapping is two-band --- score \(\leq 58\) \emph{Fragile}, \(> 58\)
\emph{Not fragile} --- with \textbf{no Robust tier}, because a single binary comparison
does not carry enough information to support a third band at the required
error rates. It applies only within the study's eligibility profile: per-arm
\(n\) in \([25, 200]\), allocation ratio in \([0.8, 1.25]\), control-arm event rate
in \([0.08, 0.55]\) with at least three events and three non-events. Outside
that profile, or under default weights, the label is suppressed.

\subsection{ANCOVA: a named negative result}\label{ancova-a-named-negative-result}

Covariate-adjusted analysis is an important calibration target because many
confirmatory trials adjust for baseline. Three pre-registered studies were
conducted; all closed without a usable label under their frozen designs.

\textbf{ANCOVA Study 1} evaluated a three-band mapping with default weights. No
threshold met the pre-specified training requirements, so held-out data was
not opened. Within its canonical Gaussian scenarios the composite score was
nearly monotone in the \(p\)-value (Spearman correlation between its bootstrap
component and \(-\log_{10}p\) was approximately 0.99). The result establishes
empirical infeasibility under that design; it does not prove that every future
ANCOVA score or scenario definition must fail.

\textbf{ANCOVA Study 2} evaluated a two-band mapping with jackknife-light weights.
Its sealed pilot showed score-location separation (\(\Delta=24.4\), overlap
0.034, AUC 0.892), but production training found no feasible cutoff: the best
clear-effect identification at a false-reassurance-safe threshold was 0.554,
below the pre-specified 0.70 requirement. This discrepancy shows why pilot
location metrics alone are insufficient for deciding whether a categorical
claim is feasible.

The \textbf{assumption-violation study} then asked whether the score added
discrimination beyond \(p\) for violations represented in a frozen scenario
suite. It completed 2,100 analyses with zero failures. The pooled difference
\(\mathrm{AUC}_\mathrm{score}-\mathrm{AUC}_p\) was 0.0053, with a
scenario-cluster bootstrap 95\% interval \([-0.111,0.129]\); the added
discrimination was not confirmed.

The ANCOVA registry rows remain \texttt{uncalibrated}, and labels stay suppressed.
Audit identifiers, hashes, and raw decision strings are retained in the
study-specific \texttt{published/} directories rather than in reader-facing prose.

\section{Examples}\label{examples}

\subsection{Welch case study}\label{welch-case-study}

The package ships a synthetic phase II analgesic trial: 55 patients, change
from baseline in pain score at week 12. Results below are loaded from a
committed artifact generated with the package-default \texttt{n\_boot\ =\ 1000} and
\texttt{seed\ =\ 123}. Welch is numeric-only after the prospective negative result.

\begin{verbatim}
res <- robustness_analysis(pain_treatment, pain_placebo, test_type = "t.test",
                           n_boot = 1000, seed = 123, interpret = TRUE)
\end{verbatim}

The Welch test gives a mean difference of -11.40
points, \(p = 0.0018\). The robustness analysis
returns a numeric score of \textbf{71.9/100} with no
categorical label, built
from: jackknife stability 100.0\% with
0 influential observations flagged; a
worst-case fragility index of \(k = 6\)
(10.9\% of the sample), with \(p\) rising to
0.0602 after the sixth deletion; and bootstrap
same-decision rate 88.7\%.

Figure \ref{fig:casefig} shows the two diagnostic displays produced by
\texttt{plot()}. The upper panel makes the contrast of the Methods section concrete:
the worst-case trajectory crosses \(\alpha\) well before the extreme-value
trajectory does.

\begin{figure}
\includegraphics[width=1\linewidth,alt={Two Welch case-study diagnostics. The upper plot compares p-value trajectories under adversarial and outlier-ranked deletion; the lower histogram shows bootstrap p-values relative to the observed p-value.}]{stabilitest_files/figure-latex/casefig-1} \caption{Diagnostic plots for the Welch case study. Top: p-value trajectories under worst-case (adversarial) and extreme-value (outlier-ranked) removal, with the dashed line at alpha = 0.05. Bottom: bootstrap p-value distribution, red line at the observed p-value.}\label{fig:casefig}
\end{figure}

For a reader, the actionable output is not the score but the six named
patients --- led by one extreme responder --- whose charts warrant source-data
review. The observed component pattern resembles the study's clear-effect
scenarios more than its null scenarios, but a sensitivity diagnostic does not
establish which data-generating truth produced an individual dataset.

\subsection{Other engines}\label{other-engines}

Short live examples across the remaining engines. An ANCOVA term on the
packaged illustration dataset:

\begin{verbatim}
res_lm <- robustness_lm(change ~ arm + baseline_pain, data = pain_ancova_trial,
                        term = "arm", n_boot = 300, seed = 123)
res_lm$robustness_metrics[, c("worstcase_fragility_k", "worstcase_fragility_pct",
                              "bootstrap_reproducibility", "overall_robustness")]
\end{verbatim}

\begin{verbatim}
#> # A tibble: 1 x 4
#>   worstcase_fragility_k worstcase_fragility_pct bootstrap_reproducibility
#>                   <int>                   <dbl>                     <dbl>
#> 1                     7                    8.75                      83.3
#> # i 1 more variable: overall_robustness <dbl>
\end{verbatim}

The score is numeric-only: 67.9
with no label, exactly as the negative result above requires. A TOST
equivalence analysis behaves the same way:

\begin{verbatim}
res_tost <- robustness_tost(pain_treatment, pain_placebo, type = "equivalence",
                            endpoint = "mean", delta_L = -20, delta_U = 20,
                            n_boot = 300, seed = 123)
c(fragility_k = res_tost$robustness_metrics$worstcase_fragility_k,
  score = round(res_tost$robustness_metrics$overall_robustness, 1))
\end{verbatim}

\begin{verbatim}
#> fragility_k       score 
#>         4.0        65.4
\end{verbatim}

A Fisher analysis inside the calibrated profile, under the calibrated weights,
does receive a label:

\begin{verbatim}
set.seed(2026)
ctrl <- rbinom(60, 1, 0.15); trt <- rbinom(60, 1, 0.40)
res_f <- robustness_analysis(trt, ctrl, test_type = "fisher", n_boot = 1000,
                             seed = 123,
                             weights = c(jackknife = 0, fragility = 0.5,
                                         bootstrap = 0.5))
c(score = round(res_f$robustness_metrics$overall_robustness, 1),
  label = res_f$robustness_interpretation)
\end{verbatim}

\begin{verbatim}
#>         score         label 
#>        "67.6" "Not fragile"
\end{verbatim}

The same data under default weights returns identical component metrics but a
different composite (75.7)
and a suppressed label --- the registry requires the calibrated weight vector,
not merely the calibrated method.

\FloatBarrier

\section{Performance}\label{performance}

Cost is dominated by the greedy search, which performs \(O(N \cdot k_{\max})\)
test evaluations. For total \(N=55\), \(k_{\max}=16\) and the greedy candidate
refits total \(\sum_{j=40}^{55}j=760\). A full default run uses approximately
\(1+55+760+16+1000=1832\) test evaluations before overhead: the original test,
jackknife, greedy search, descriptive extreme-value path, and bootstrap.
Closed-form tests are therefore sub-second; model engines pay for repeated
refits.

\begin{longtable}[t]{lrr}
\caption{\label{tab:timingtable}Wall-clock seconds for a single analysis with 1000 bootstrap resamples, stabilitest 0.6.0 on an Apple aarch64 laptop. The Cox benchmark uses simulated survival data generated in the timing script.}\\
\toprule
Engine & n = 50 per arm (s) & n = 150 per arm (s)\\
\midrule
robustness\_analysis (Welch) & 0.53 & 1.43\\
robustness\_lm (ANCOVA) & 0.98 & 4.22\\
robustness\_glm (binomial) & 1.62 & 8.03\\
robustness\_surv (Cox) & 1.64 & 6.17\\
robustness\_tost (mean equivalence) & 0.77 & 1.48\\
\bottomrule
\end{longtable}

The practical implication is that interactive use is comfortable at trial
scale, and that a large calibration study --- thousands of replications ---
belongs in a batch script, which is how the studies reported here were run.

\section{Testing and quality assurance}\label{testing-and-quality-assurance}

The package carries 1011 passing \texttt{testthat} assertions across
7 test files, covering the analysis engines, edge cases (degenerate groups, zero
variance, complete separation, failed refits), the packaged datasets, the
calibration registry's resolution logic, and the calibration documentation
itself.

That last one is unusual enough to describe. A source-tree documentation audit,
run as part of the test suite,
asserts that the README, NEWS, vignettes, manuscript, and every study SAP
still state the policy invariants: that uncalibrated labels are suppressed,
that the historical Welch mapping is inactive, and that each frozen acceptance
criterion remains present in its own statistical analysis plan. The intent is
to make it mechanically hard to quietly weaken a published calibration claim
through documentation drift.

Each calibration study additionally publishes compact decision artifacts in
its own \texttt{published/} directory: occupancy and failure tables, power
verification, candidate diagnostics, a registry row, manifests, and a hash
ledger. Candidate hashes are recorded before held-out data is opened, so
the sequencing of any decision is auditable after the fact.

\section{Discussion and limitations}\label{discussion-and-limitations}

The framework separates three questions usually blurred together, and its
sharpest practical lesson is the gap between adversarial and symmetric
removal: in the descriptive stress simulation, adversarial deletion and
cumulative grand-mean extreme-value deletion produced materially different
fragility summaries. That comparison does not benchmark Cook's distance,
DFBETAs, \CRANpkg{car}, or \CRANpkg{influence.ME}.

Several limitations deserve statement rather than burial.

\textbf{The score is largely a monotone transform of the \(p\)-value.} In clean
parametric settings the composite score re-expresses the evidence the
\(p\)-value already carries; the ANCOVA program measured this directly. This is
a scope condition, not a defect, but it bounds what the score can be used for.
Its added value lies in (i) identifying the specific observations carrying the
conclusion, (ii) the worst-case deletion bound, and (iii) behaviour under
contamination, where the \(p\)-value alone is silent --- see the null-with-
contamination rows of Table \ref{tab:simtable}. The assumption-violation study
tested whether the score added discrimination beyond \(p\) under its frozen
violation suite and did not confirm such an improvement.

\textbf{The weights are a convention.} They are explicit, pre-specifiable, and
demonstrably not load-bearing for discrimination (the Fisher weight study
found AUC essentially unchanged across weightings), but they are not derived
from theory. The score is a communication device; the components are the
substance.

\textbf{The greedy index is an upper bound.} Exact minimal overturning subsets are
combinatorially hard. The algorithm returns a sufficient overturning set when
it finds one, not a certificate of minimality. Certified minimal sets --- via
integer-programming formulations --- remain future work.

\textbf{Binary significance framing.} Inherited from the fragility concept itself
and criticised for good reason \citep{potter2020}; the package's response is to
report the trajectory alongside the index and to frame the index as an upper
bound, wherever the index is defined rather than only in a limitations
section.

\textbf{Independence is assumed.} Clustered and longitudinal designs are not
supported; for those, \texttt{influence.ME} addresses the grouping-factor question.

\textbf{Calibration coverage is deliberately narrow.} Only the eligible Fisher
exact-test configuration under jackknife-light weights carries a label.
Welch, paired, rank, Brunner--Munzel, non-canonical binary, model-based, Cox,
and TOST analyses return numeric output with a suppressed label pending an
independent positive calibration. Users who want a verdict where none is
licensed are, by design, not given one.

\section{Summary}\label{summary}

\CRANpkg{stabilitest} asks how much of the data a conclusion actually rests on,
and answers with component metrics that name the observations responsible. Its
methodological position is that a categorical robustness verdict is a claim
requiring validation for the specific analysis at hand, and that the honest
response to a failed validation is silence plus a published negative result.
The package is available from
\href{https://github.com/ma-brain/stabilitest}{github.com/ma-brain/stabilitest}.

\bibliography{RJreferences.bib}

\address{%
Marius Ardelean\\
Independent researcher\\%
\\
%
\url{https://github.com/ma-brain/stabilitest}\\%
%
\href{mailto:marally@gmail.com}{\nolinkurl{marally@gmail.com}}%
}
